\documentclass{article}
\usepackage{tikz}
\usetikzlibrary{arrows.meta,positioning,calc,shapes.geometric}
\pagestyle{empty}
\begin{document}
\begin{tikzpicture}

\draw[->] (0,-1.2) -- (0,5.6) node[above, font=\scriptsize] {potencial (mV)};
\draw[->] (0,0) -- (9.6,0) node[right, font=\scriptsize] {tiempo};
\draw[dashed] (0,-0.35) -- (9,-0.35) node[right, font=\scriptsize] {reposo $\approx -70$};
\draw[dashed] (0,0.25) -- (9,0.25) node[right, font=\scriptsize] {umbral $\approx -55$};
\draw[thick, blue] plot[smooth] coordinates {(0,-0.35) (1,-0.35) (1.7,0.25) (2.1,2.2) (2.5,4.4) (2.9,4.7) (3.3,3.2) (3.8,0.5) (4.4,-0.9) (5.0,-1.1) (5.8,-0.7) (6.8,-0.4) (8.5,-0.35)};
\node[font=\scriptsize] at (2.2,5.3) {1 ascenso}; \node[font=\scriptsize] at (3.6,5.0) {2 pico};
\node[font=\scriptsize] at (5.1,2.4) {3 descenso}; \node[font=\scriptsize] at (6.2,-1.5) {4 sobrepasamiento};
\node[font=\scriptsize, align=center] at (7.5,-2.3) {5 periodo refractario (solapado con 3, 4)};
\node[font=\scriptsize, align=center] at (0.9,1.0) {estímulo};
\draw[->] (0.9,0.6) -- (1.0,-0.25);

\end{tikzpicture}
\end{document}
